\section{Lean formalization}\label{app:formalization}
\subsection*{Trust base and methodology}

The Lean development for this paper is mathlib-free.  It imports no
external mathematics library, so no mathlib axioms enter the dependency
closure.  Every declaration registered as a theorem was audited with
\texttt{\#print axioms}.  The only axioms that occur are Lean's standard
foundational ones: \texttt{propext}, \texttt{Quot.sound}, and
\texttt{Classical.choice}, with \texttt{funext} and the \texttt{choice}
recursor appearing as their standard consequences in the audited closures.
There are no project-local axioms behind any theorem-marked result.  Thus
the trust base for the formalized theorem claims is precisely Lean's kernel
plus classical logic, and nothing project-specific.

This limited trust base is possible because the formalized content is
finite and typed.  The adequacy collapses and no-go residuals are not
mechanized as analytic estimates, operator-theoretic assertions, or
motivic constructions.  They are finite vector and matrix computations:
the source and target covariance blocks are explicit finite matrices, the
full-source collapses reduce to identity-matrix action, and the scalar
shadow residuals are computed from rank-\(1\) projectors over \(\Q\).
Each Schur-complement residual is therefore a decidable calculation in a
small finite algebraic model.  This is the reason a mathlib-free kernel can
check the main adequacy and no-go statements.

The surrounding arithmetic geometry is treated separately from those finite
calculations.  The Lean development checks the typed identities appearing
in the apparatus: for example, that an identity source kills the Schur
residual, that the relevant projector traces have the stated values, and
that the finite noninjectivity witnesses have the claimed public shadows.
It does not certify Gross--Zagier, Iwasawa main conjecture inputs,
Cassels--Tate comparison theory, Bloch--Kato local comparison maps, or the
motivic descents named in the body.  Such arithmetic content enters the
paper either as cited external input, as a recognition-source carrier, or
as an explicitly stated hypothesis.

The statements-of-record inventory for this paper contains \(41\) items.
Of these, \(19\) are theorem-grade Lean results with substantive finite
proofs, \(11\) are typed definition mirrors, and \(11\) are not mechanized:
the five recognition-source transport obligations, the four support-only
remarks, the warning, and the scope nonclaim.  Those non-mechanized items
are paper-level statements carried as motivation, scope, or hypotheses,
not Lean declarations.  Two mechanized items, the \(\vsrc\) algebra
definition and the Stage-I \(\vsrc\) consistency theorem, are checked over
a narrowed structural representation; the following subsections record the
per-label coverage, declaration map, and standing narrowings in detail.
\subsection*{Wording-discipline table}

Each theorem in the body carries exactly one mechanization footnote, and
the wording of that footnote is canonical rather than editorial.  The
wording records the fidelity class of the formal counterpart: a faithful
Lean theorem, a Lean theorem over a narrowed representation, a typed
definition mirror, or an unmechanized paper-level statement.  Narrowing and
non-mechanization are therefore visible in the disclosure convention rather
than hidden in prose.

\begin{center}
\begin{tabularx}{\textwidth}{@{}>{\raggedright\arraybackslash}p{0.31\textwidth}X@{}}
\toprule
Statement class & Footnote convention \\
\midrule
Theorem, faithful &
Verified in Lean as \texttt{[decl]}; see App~\ref{app:formalization}. \\
\addlinespace
Theorem, narrowed representation &
Verified in Lean as \texttt{[decl]} over a narrowed representation
(e.g.\ the \(2^r\)-coordinate space); the narrowing is recorded below.  The
narrowing qualifier is mandatory; never bare ``verified in Lean''. \\
\addlinespace
Definition, narrowed representation &
Realized in Lean as \texttt{[decl]} over a narrowed representation; the
narrowing is recorded below. \\
\addlinespace
Definition, faithful (typed mirror) &
No footnote; the declaration appears only in the per-label coverage table
below. \\
\addlinespace
Obligation / remark / warning / non-claim (not mechanized) &
No footnote and no Lean identifier in the body; recorded as a paper-level
statement only. \\
\bottomrule
\end{tabularx}
\end{center}
\subsection*{Standing narrowings}

There are two standing qualifications behind the coverage claims in this
appendix.  Both are part of the formalization contract: they delimit what
the finite Lean development checks, and they keep the surrounding
arithmetic content visible rather than implicit.

First, two \(\vsrc\) items are checked over a narrowed structural
representation: the \(\vsrc\) algebra definition and the Stage-I
\(\vsrc\) consistency theorem.  On the \(2^r\)-coordinate model, the
determinant-shaped content of a rank-\(r\) source is represented, and the
low-rank exterior relations actually used in the paper are checked.  In
rank \(2\), this includes the wedge relation
\[
  \Jone\Jtwo=\Jtwelve
\]
together with the accompanying nilpotence and anticommutation relations.
The dimension count and the calibration of the determinant shape are
faithful to the intended exterior-algebra source.

What is not claimed is equally important.  The formalization is not a
complete general exterior-algebra development: it does not provide general
\(\bigwedge^k\) functoriality, nor a full axiomatization of exterior
algebras in arbitrary rank.  It is also not a descent from \(\vsrc\) to
actual motivic cycles.  The Lean check certifies the finite structural
relations on the coordinate representation, not a general motivic
construction.  This is why those two body items carry the ``narrowed
representation'' footnote convention rather than a bare ``verified in
Lean'' disclosure, and why the \(\vsrc\) statements in
\Cref{sec:vsrc} are calibration and consistency results rather than
constructions of motivic sources.

Second, the five per-column transports are recognition sources, not
derived theorems and not Lean declarations.  These are the
height/regulator, analytic/period, finite-source, \(p\)-adic, and
determinant-assembly comparison transports appearing in
\Cref{sec:height-regulator,sec:analytic-period,sec:finite-source,sec:p-adic,sec:det-assembly}.
Each column's Schur-adequacy collapse is proved at the finite typed level
once the column's native source is admitted.  The corresponding transport
is precisely the named arithmetic comparison content that admits that
source into the Bloch--Kato determinant-line target.  It is supplied or
imported externally and then carried as a hypothesis.

Thus the apparatus does not prove these five transports.  It records them
as obligations and proves the per-column collapses conditional on the
native source data that those transports are meant to recognize.  In the
per-label coverage table below they appear as not-mechanized obligations.
Neither standing narrowing enlarges the Lean trust base: no new axioms are
introduced; the qualifications bound the finite mechanization claims, with
the remaining arithmetic content carried explicitly.
\begin{table}[tbp]
\centering
\small
\caption{Standing representation notes for the apparatus formalization.}
\label{tab:representation-notes}
\begin{tabularx}{\textwidth}{@{}>{\raggedright\arraybackslash}p{0.24\textwidth}
>{\raggedright\arraybackslash}p{0.28\textwidth}
>{\raggedright\arraybackslash}X@{}}
\toprule
Representation issue & Where it appears & Disclosure effect \\
\midrule
\(\vsrc\) structural encoding &
\(\Cref{sec:vsrc}\), App~\ref{app:formalization} &
The rank-\(2\) exterior relations are checked on the \(2^r\)-coordinate
model.  The paper claims calibration and consistency, not a full exterior
algebra or motivic descent formalization. \\
\addlinespace
\(\Q\) rank-\(1\) projector traces &
Finite-source and no-go statements &
Scalar-shadow residuals are finite rational projector computations; the
Lean result is the finite trace calculation, not an arithmetic
classification theorem. \\
\addlinespace
Identity-matrix action &
Schur-collapse statements &
Full-source adequacy reduces to \(I\)-action on an explicit finite carrier.
The external transport that admits the native source remains a carried
recognition obligation. \\
\bottomrule
\end{tabularx}
\end{table}

\subsection*{Per-label coverage}

The following table lists every statement of record in document order,
with its environment kind, its mechanization fidelity, and its body
location.  The paper labels in the first column elide the
\texttt{apparatus:} prefix.  Faithful and narrowed rows have Lean
declarations, mapped in the next subsection; not-mechanized rows do not.

\begin{small}
\begin{longtable}{@{}>{\ttfamily}p{0.45\textwidth} l l l@{}}
\caption{Lean coverage of the labelled statements in the inventory, in document order (41 rows).  Labels omit the \texttt{apparatus:} infix.}\label{tab:formalization-coverage}\\
\toprule \normalfont Label & Kind & Lean coverage & Section \\ \midrule \endfirsthead
\multicolumn{4}{@{}l}{\footnotesize\itshape continued from previous page}\\ \toprule \normalfont Label & Kind & Lean coverage & Section \\ \midrule \endhead
\midrule \multicolumn{4}{r@{}}{\footnotesize\itshape continued on next page}\\ \endfoot
\bottomrule \endlastfoot
def:bk-component-map & \normalfont definition & \normalfont Lean definition & \normalfont \Cref{sec:decomposition} \\
def:bridge-defect-equation & \normalfont definition & \normalfont Lean definition & \normalfont \Cref{sec:decomposition} \\
thm:component-killing & \normalfont theorem & \normalfont Lean proves & \normalfont \Cref{sec:decomposition} \\
thm:finite-scalar-public-shadow & \normalfont theorem & \normalfont model version & \normalfont \Cref{sec:decomposition} \\
nonclaim:decomposition-scope & \normalfont nonclaim & \normalfont not formalized & \normalfont \Cref{sec:decomposition} \\
def:height-regulator-map & \normalfont definition & \normalfont Lean definition & \normalfont \Cref{sec:height-regulator} \\
thm:height-schur-collapse & \normalfont theorem & \normalfont model version & \normalfont \Cref{sec:height-regulator} \\
obl:rho-ht-reg-transport & \normalfont obligation & \normalfont not formalized & \normalfont \Cref{sec:height-regulator} \\
warn:571a1-carrier & \normalfont warning & \normalfont not formalized & \normalfont \Cref{sec:height-regulator} \\
def:analytic-period-map & \normalfont definition & \normalfont Lean definition & \normalfont \Cref{sec:analytic-period} \\
thm:analytic-period-schur-collapse & \normalfont theorem & \normalfont Lean proves & \normalfont \Cref{sec:analytic-period} \\
obl:rho-an-period-transport & \normalfont obligation & \normalfont not formalized & \normalfont \Cref{sec:analytic-period} \\
def:finite-source-map & \normalfont definition & \normalfont Lean definition & \normalfont \Cref{sec:finite-source} \\
thm:finite-source-schur-collapse & \normalfont theorem & \normalfont Lean proves & \normalfont \Cref{sec:finite-source} \\
thm:scalar-finite-shadow & \normalfont theorem & \normalfont Lean proves & \normalfont \Cref{sec:finite-source} \\
thm:dim-sha-shadow & \normalfont theorem & \normalfont model version & \normalfont \Cref{sec:finite-source} \\
obl:rho-finite-transport & \normalfont obligation & \normalfont not formalized & \normalfont \Cref{sec:finite-source} \\
def:p-adic-map & \normalfont definition & \normalfont Lean definition & \normalfont \Cref{sec:p-adic} \\
thm:ordinary-schur-collapse & \normalfont theorem & \normalfont model version & \normalfont \Cref{sec:p-adic} \\
thm:signed-schur-collapse & \normalfont theorem & \normalfont model version & \normalfont \Cref{sec:p-adic} \\
obl:rho-p-transport & \normalfont obligation & \normalfont not formalized & \normalfont \Cref{sec:p-adic} \\
def:det-assembly-map & \normalfont definition & \normalfont Lean definition & \normalfont \Cref{sec:det-assembly} \\
thm:cascade-completion-signature & \normalfont theorem & \normalfont Lean proves & \normalfont \Cref{sec:det-assembly} \\
obl:rho-det-assembly & \normalfont obligation & \normalfont not formalized & \normalfont \Cref{sec:det-assembly} \\
rmk:five-column-literature-gap & \normalfont remark & \normalfont not formalized & \normalfont \Cref{sec:comparison} \\
def:five-column-tensor & \normalfont definition & \normalfont Lean definition & \normalfont \Cref{sec:comparison} \\
thm:five-col-bk-comparison & \normalfont theorem & \normalfont Lean proves & \normalfont \Cref{sec:comparison} \\
def:rank-two-kummer-source & \normalfont definition & \normalfont Lean definition & \normalfont \Cref{sec:higher-rank-no-go} \\
thm:rank-two-height-schur-collapse & \normalfont theorem & \normalfont model version & \normalfont \Cref{sec:higher-rank-no-go} \\
thm:gram-matrix-shadow & \normalfont theorem & \normalfont model version & \normalfont \Cref{sec:higher-rank-no-go} \\
thm:finite-window-no-go & \normalfont theorem & \normalfont model version & \normalfont \Cref{sec:global-audit-no-go} \\
thm:finite-prime-cover-no-go & \normalfont theorem & \normalfont Lean proves & \normalfont \Cref{sec:support-prime-no-go} \\
def:support-prime-truncation-residual & \normalfont definition & \normalfont Lean definition & \normalfont \Cref{sec:support-prime-no-go} \\
rmk:support-prime-atlas & \normalfont remark & \normalfont not formalized & \normalfont \Cref{sec:support-prime-no-go} \\
def:vsrc-algebra & \normalfont definition & \normalfont Lean definition & \normalfont \Cref{sec:vsrc} \\
thm:vsrc-stage-i-consistency & \normalfont theorem & \normalfont model version & \normalfont \Cref{sec:vsrc} \\
rmk:vsrc-no-licensed-descent & \normalfont remark & \normalfont not formalized & \normalfont \Cref{sec:vsrc} \\
thm:vsrc-stage-ii-gz-partial & \normalfont theorem & \normalfont model version & \normalfont \Cref{sec:vsrc} \\
thm:vsrc-stage-iii-translation & \normalfont theorem & \normalfont conditional schema & \normalfont \Cref{sec:vsrc} \\
rmk:vsrc-rank-r-stronger & \normalfont remark & \normalfont not formalized & \normalfont \Cref{sec:vsrc} \\
thm:kappa-r-normalization-equivalence & \normalfont theorem & \normalfont model version & \normalfont \Cref{sec:kappa-normalization} \\
\end{longtable}
\end{small}

\subsection*{Lean declaration map}

The \(30\) mechanized statements, consisting of the \(11\) typed
definitions and \(19\) theorems, each correspond to a Lean declaration
listed here for citation and auditing.  Paper labels elide the
\texttt{apparatus:} prefix, and declarations elide the
\texttt{SixBirdsBSD.Apparatus.} namespace.  The not-mechanized
obligations, remarks, warning, and nonclaim have no declaration and are
therefore omitted; the full coverage table, including those rows, is
\Cref{tab:formalization-coverage}.

\begin{small}
\begin{longtable}{@{}>{\ttfamily\raggedright\arraybackslash}p{0.44\textwidth} >{\ttfamily\raggedright\arraybackslash}p{0.50\textwidth}@{}}
\caption{Lean declaration map for the 30 mechanized apparatus statements (definitions and theorems). Paper labels elide the \texttt{apparatus:} prefix; declarations elide the \texttt{SixBirdsBSD.Apparatus.} namespace.}\label{tab:appd-declmap}\\
\toprule \normalfont Paper label & \normalfont Lean declaration \\ \midrule
\endfirsthead
\multicolumn{2}{@{}l}{\footnotesize\itshape continued from previous page}\\
\toprule \normalfont Paper label & \normalfont Lean declaration \\ \midrule
\endhead
\midrule \multicolumn{2}{r@{}}{\footnotesize\itshape continued on next page}\\
\endfoot
\bottomrule
\endlastfoot
def:\allowbreak bk-\allowbreak component-\allowbreak map & Decomposition.\allowbreak bk\allowbreak Component\allowbreak Map \\
def:\allowbreak bridge-\allowbreak defect-\allowbreak equation & Decomposition.\allowbreak bridge\allowbreak Defect\allowbreak Equation \\
thm:\allowbreak component-\allowbreak killing & Decomposition.\allowbreak component\allowbreak Killing \\
thm:\allowbreak finite-\allowbreak scalar-\allowbreak public-\allowbreak shadow & Decomposition.\allowbreak finite\allowbreak Scalar\allowbreak Public\allowbreak Shadow \\
def:\allowbreak height-\allowbreak regulator-\allowbreak map & Height\allowbreak Regulator.\allowbreak height\allowbreak Regulator\allowbreak Map \\
thm:\allowbreak height-\allowbreak schur-\allowbreak collapse & Height\allowbreak Regulator.\allowbreak height\allowbreak Schur\allowbreak Collapse \\
def:\allowbreak analytic-\allowbreak period-\allowbreak map & Analytic\allowbreak Period.\allowbreak analytic\allowbreak Period\allowbreak Map \\
thm:\allowbreak analytic-\allowbreak period-\allowbreak schur-\allowbreak collapse & Analytic\allowbreak Period.\allowbreak analytic\allowbreak Period\allowbreak Schur\allowbreak Collapse \\
def:\allowbreak finite-\allowbreak source-\allowbreak map & Finite\allowbreak Source.\allowbreak finite\allowbreak Source\allowbreak Map \\
thm:\allowbreak finite-\allowbreak source-\allowbreak schur-\allowbreak collapse & Finite\allowbreak Source.\allowbreak finite\allowbreak Source\allowbreak Schur\allowbreak Collapse \\
thm:\allowbreak scalar-\allowbreak finite-\allowbreak shadow & Finite\allowbreak Source.\allowbreak scalar\allowbreak Finite\allowbreak Shadow \\
thm:\allowbreak dim-\allowbreak sha-\allowbreak shadow & Finite\allowbreak Source.\allowbreak dim\allowbreak Sha\allowbreak Shadow \\
def:\allowbreak p-\allowbreak adic-\allowbreak map & PAdic.\allowbreak p\allowbreak Adic\allowbreak Map \\
thm:\allowbreak ordinary-\allowbreak schur-\allowbreak collapse & PAdic.\allowbreak ordinary\allowbreak Schur\allowbreak Collapse \\
thm:\allowbreak signed-\allowbreak schur-\allowbreak collapse & PAdic.\allowbreak signed\allowbreak Schur\allowbreak Collapse \\
def:\allowbreak det-\allowbreak assembly-\allowbreak map & Det\allowbreak Assembly.\allowbreak det\allowbreak Assembly\allowbreak Map \\
thm:\allowbreak cascade-\allowbreak completion-\allowbreak signature & Det\allowbreak Assembly.\allowbreak cascade\allowbreak Completion\allowbreak Signature \\
def:\allowbreak five-\allowbreak column-\allowbreak tensor & Comparison.\allowbreak five\allowbreak Column\allowbreak Tensor \\
thm:\allowbreak five-\allowbreak col-\allowbreak bk-\allowbreak comparison & Comparison.\allowbreak five\allowbreak Col\allowbreak Bk\allowbreak Comparison \\
def:\allowbreak rank-\allowbreak two-\allowbreak kummer-\allowbreak source & Higher\allowbreak Rank\allowbreak No\allowbreak Go.\allowbreak rank\allowbreak Two\allowbreak Kummer\allowbreak Source \\
thm:\allowbreak rank-\allowbreak two-\allowbreak height-\allowbreak schur-\allowbreak collapse & Higher\allowbreak Rank\allowbreak No\allowbreak Go.\allowbreak rank\allowbreak Two\allowbreak Height\allowbreak Schur\allowbreak Collapse \\
thm:\allowbreak gram-\allowbreak matrix-\allowbreak shadow & Higher\allowbreak Rank\allowbreak No\allowbreak Go.\allowbreak gram\allowbreak Matrix\allowbreak Shadow \\
thm:\allowbreak finite-\allowbreak window-\allowbreak no-\allowbreak go & Global\allowbreak Audit\allowbreak No\allowbreak Go.\allowbreak finite\allowbreak Window\allowbreak No\allowbreak Go \\
thm:\allowbreak finite-\allowbreak prime-\allowbreak cover-\allowbreak no-\allowbreak go & Support\allowbreak Prime\allowbreak No\allowbreak Go.\allowbreak finite\allowbreak Prime\allowbreak Cover\allowbreak No\allowbreak Go \\
def:\allowbreak support-\allowbreak prime-\allowbreak truncation-\allowbreak residual & Support\allowbreak Prime\allowbreak No\allowbreak Go.\allowbreak support\allowbreak Prime\allowbreak Truncation\allowbreak Residual \\
def:\allowbreak vsrc-\allowbreak algebra & Vsrc.\allowbreak vsrc\allowbreak Algebra \\
thm:\allowbreak vsrc-\allowbreak stage-\allowbreak i-\allowbreak consistency & Vsrc.\allowbreak vsrc\allowbreak Stage\allowbreak IConsistency \\
thm:\allowbreak vsrc-\allowbreak stage-\allowbreak ii-\allowbreak gz-\allowbreak partial & Vsrc.\allowbreak vsrc\allowbreak Stage\allowbreak IIGz\allowbreak Partial \\
thm:\allowbreak vsrc-\allowbreak stage-\allowbreak iii-\allowbreak translation & Vsrc.\allowbreak vsrc\allowbreak Stage\allowbreak IIITranslation \\
thm:\allowbreak kappa-\allowbreak r-\allowbreak normalization-\allowbreak equivalence & Kappa\allowbreak Normalization.\allowbreak kappa\allowbreak RNormalization\allowbreak Equivalence \\
\end{longtable}
\end{small}
